% =============================================================================
%  main_paper_c.tex
%  Journal Paper C — IEEE Transactions on Power Systems
%
%  Title: Self-Adaptive Model-Based Protection: A Unified Framework for
%         Real-Time Inverse Estimation Across Overcurrent, Transformer,
%         Line, and Busbar Differential Protection
%
%  Compile:
%    pdflatex main_paper_c && bibtex main_paper_c &&
%    pdflatex main_paper_c && pdflatex main_paper_c
% =============================================================================

\documentclass[journal,twoside,10pt]{IEEEtran}

% ---- Essential packages -----------------------------------------------------
\usepackage[T1]{fontenc}
\usepackage[utf8]{inputenc}
\usepackage{mathptmx}
\usepackage{microtype}

% ---- Mathematics ------------------------------------------------------------
\usepackage{amsmath,amssymb,mathtools}
\usepackage{amsthm}
\usepackage{bm}
\usepackage{siunitx}
\sisetup{detect-all,per-mode=symbol,range-phrase=\,--\,}

% ---- Figures and tables -----------------------------------------------------
\usepackage{graphicx}
\usepackage{subcaption}
\usepackage{booktabs}
\usepackage{array}
\usepackage{multirow}
\usepackage{float}
\usepackage{rotating}

% ---- References -------------------------------------------------------------
\usepackage[colorlinks=false]{hyperref}
\usepackage{xcolor}
\bibliographystyle{IEEEtran}

% ---- Algorithm --------------------------------------------------------------
\usepackage{algorithm}
\usepackage{algpseudocode}
\renewcommand{\algorithmicrequire}{\textbf{Input:}}
\renewcommand{\algorithmicensure}{\textbf{Output:}}

% ---- Theorem environments ---------------------------------------------------
\newtheorem{theorem}{Theorem}
\newtheorem{proposition}{Proposition}
\newtheorem{remark}{Remark}
\newtheorem{corollary}{Corollary}

% ---- Custom macros ----------------------------------------------------------
\newcommand{\vect}[1]{\bm{#1}}
\newcommand{\mat}[1]{\mathbf{#1}}
\newcommand{\norm}[1]{\left\|#1\right\|}
\newcommand{\pu}{\,\text{pu}}
\newcommand{\mss}{\,\text{ms}}
\newcommand{\sss}{\,\text{s}}
\newcommand{\SAMBP}{\textsc{sambp}}

% =============================================================================
\begin{document}
% =============================================================================

\title{Self-Adaptive Model-Based Protection: A Unified Framework
       for Real-Time Inverse Estimation Across Overcurrent,
       Transformer, Line, and Busbar Differential Protection}

\author{Anoop~V.~Eluvathingal,~\IEEEmembership{Student Member,~IEEE,}
        and~K.~Shanti~Swarup,~\IEEEmembership{Senior Member,~IEEE}%
\thanks{Manuscript received \today.
        This work is part of the PhD research programme at the Department of
        Electrical Engineering, Indian Institute of Technology Madras,
        Chennai~600~036, India.}%
\thanks{A.~V.~Eluvathingal and K.~S.~Swarup are with the Department of
        Electrical Engineering, IIT Madras, Chennai~600~036, India
        (e-mail: ee21d017@smail.iitm.ac.in; swarup@ee.iitm.ac.in).}}

\markboth{IEEE Transactions on Power Systems,~Vol.~XX, No.~XX, XXXX~2026}%
{Eluvathingal \& Swarup: SAMBP — Unified Inverse Estimation Framework}

\maketitle

% =============================================================================
\begin{abstract}
Protection relays across overcurrent (OC), transformer differential (87T),
line current differential (87L), and busbar differential (87B) functions
use fixed-commissioning settings that are blind to operating conditions,
fault impedance, and network topology changes.
This paper presents \textsc{sambp} (Self-Adaptive Model-Based Protection),
a unified five-layer architecture applied uniformly across all four
protection functions: a conventional baseline relay (Layer~1), a
physics-constrained reduced-parameter zone model (Layer~2), a two-pass
Levenberg--Marquardt estimator (Layer~3), a composite confidence gate
(Layer~4), and a function-specific fallback or coordination layer (Layer~5).
The key architectural invariant is a physics-parameter-reduction discipline
— each zone model eliminates one collinear or collinear-via-physics parameter
by derivation rather than fitting, reducing the Jacobian condition number
from $\mathcal{O}(10^8)$ (unconstrained) to $\kappa_n < 30$ across all
four functions.
The Stage-2 model veto ($\kappa_n < \kappa_{thresh}$ AND $f_{int} <
f_{thresh}$) provides a generalised mechanism for CT saturation rejection
(87T, 87B) and channel anomaly rejection (87L) without function-specific
hard-coding.
Validated across 28 simulation scenarios (3 OC, 7~87T, 12~87L, 6~87B),
the framework achieves correct classification in 26/28 cases; the 2
exceptions are a known architectural boundary (CT open-circuit on 87B)
and an expected no-trip case (channel-loss internal fault on 87L).
Combined compute time is below 20\,ms for all four functions, meeting
sub-cycle embedded relay requirements.
\end{abstract}

\begin{IEEEkeywords}
adaptive protection, inverse parameter estimation, Levenberg--Marquardt,
transformer differential, line differential, busbar differential,
overcurrent relay, CT saturation, confidence gate, unified framework.
\end{IEEEkeywords}

% =============================================================================
\section{Introduction}
\label{sec:intro}
% =============================================================================

Power system protection relays are engineered for worst-case fault
conditions at commissioning time.
For overcurrent (OC) protection, settings are calculated from maximum
bolted-fault current and remain fixed even as generation mix changes;
for differential protection (87T, 87L, 87B), operating and restraint
thresholds are similarly static~\cite{Blackburn2006,Anderson1999}.
This fundamental rigidity creates three classes of systematic mis-coordination
as active networks evolve:

\begin{enumerate}
  \item \textit{OC mis-calibration}: Fixed pickup settings exceed actual
        fault current under high-impedance conditions, creating security
        deficits for partial-fault scenarios~\cite{Brahma2004}.
  \item \textit{CT saturation false tripping}: Conventional harmonic-blocking
        (87T) and restrained-differential (87B) schemes can mal-operate
        under heavy through-fault CT saturation when the conventional
        threshold is exceeded before the blocking signature is
        recognised~\cite{Kasztenny2003,Schweitzer2009}.
  \item \textit{Channel anomaly} (87L): Synchronisation errors, propagation
        skew, and communication channel degradation introduce spurious
        differential current that can exceed relay thresholds during
        healthy network conditions~\cite{Adamiak2006}.
\end{enumerate}

Existing adaptive approaches address these in isolation.
Brahma and Girgis~\cite{Brahma2004} adapt OC relay settings for DG
penetration but require topology enumeration rather than direct
parameter estimation.
Saleh \textit{et al.}~\cite{Saleh2017} optimise coordination for
microgrid contingencies but assume static relay characteristics.
Sachdev and Nagpal~\cite{Sachdev1988} applied recursive least-squares
to relay measurements but without physics-derived parameter constraints,
leaving Jacobian ill-conditioning unaddressed.
No existing work applies a single estimation-based framework across all
four protection functions with a formally characterised veto mechanism.

Recent machine-learning approaches~\cite{Sheikhzadeh2022,Gao2023} demonstrate
adaptive relay classification but require large labelled datasets, offer
limited interpretability, and have no established certification path
under IEC~60255 and IEEE~C37 standards.

This paper makes four contributions:

\begin{enumerate}
  \item A \textbf{unified five-layer \textsc{sambp} architecture} applicable
        across OC, 87T, 87L, and 87B protection with a common design
        discipline and a shared LM estimator backend.
  \item A \textbf{Physics-Reduction Proposition} (§\ref{sec:comparison}):
        for any zone model with collinear or physics-derivable parameters,
        replacing fitting with derivation reduces $\kappa_n$ by at least
        one order of magnitude, enabling reliable LM convergence.
  \item A \textbf{generalised Stage-2 model veto} that detects CT
        saturation (87T, 87B) and channel anomaly (87L) as identifiability
        conditions rather than signal-processing heuristics.
  \item A \textbf{cross-function quantitative comparison}: parameter count,
        $\kappa_n$ range, event classification accuracy, and compute cost
        across all four functions (Table~\ref{tab:model_compare}).
\end{enumerate}

% =============================================================================
\section{The SAMBP Five-Layer Architecture}
\label{sec:arch}
% =============================================================================

Fig.~\ref{fig:arch} illustrates the five-layer \textsc{sambp} stack applied
uniformly to all four protection functions.

\begin{figure}[!t]
  \centering
  \includegraphics[width=\columnwidth]{figures/fig_architecture.pdf}
  \caption{Five-layer \textsc{sambp} architecture applied across four
           protection functions. Each column is one function; each row
           is one layer. All four functions share the two-pass LM
           estimator (L3) and the composite confidence gate (L4).
           Layer-2 zone models and Layer-5 fallback/coordination
           are function-specific.}
  \label{fig:arch}
\end{figure}

\subsection{Layer 1 — Conventional Baseline}
\label{sec:arch:l1}

Each protection function is backed by a fully independent conventional
relay that operates without any reference to the estimation layers.
The conventional trip signal \texttt{conv\_trip} is a \emph{necessary
condition} for any model-initiated trip: the model can only assert a trip
when the conventional relay has already asserted, preventing model-only
trips in low-current conditions where estimation uncertainty is high.
The model veto (Layer 4) is the only path by which the model can
\emph{suppress} a conventional trip.

\subsection{Layer 2 — Reduced-Parameter Zone Model}
\label{sec:arch:l2}

Each function defines a zone model $\hat{i}(\theta, t)$ with a minimal
parameter vector $\theta$ derived from the physics of the protected element.
The key design discipline is the \emph{physics-reduction step}: one
parameter that is either collinear with others or analytically determined
by network constraints is \emph{derived} rather than fitted.
Table~\ref{tab:model_compare} summarises the four zone models.

\subsection{Layer 3 — Two-Pass LM Estimator}
\label{sec:arch:l3}

All four functions share the same two-pass estimation strategy:
Pass~1 fits all parameters over a full post-fault window (100--200\,ms) to
identify the global waveform shape;
Pass~2 refines the fastest-changing parameters (decay constants, CT error)
over a shorter tail window where those parameters are best conditioned.
The SciPy \texttt{least\_squares} implementation~\cite{Virtanen2020}
is used across all four functions with function-specific Jacobian
computations.

\subsection{Layer 4 — Composite Confidence Gate and Model Veto}
\label{sec:arch:l4}

The confidence score $\gamma = w_r s_r + w_c s_c + w_b s_b + w_p s_p$
(detailed in~\cite{EluvathingalPaperA}) gates relay adaptation:
$\gamma \geq \gamma_{\min} = 0.70$ is required to accept an adaptive trip.
Beyond the forward gate, the Stage-2 model veto is activated when:
\begin{equation}
  \kappa_n < \kappa_{thresh} \;\text{ AND }\; f_{int} < f_{thresh},
  \label{eq:veto}
\end{equation}
where $f_{int} \in [0,1]$ is the derived internal-fault probability and
$\kappa_{thresh} = 30$.
When~\eqref{eq:veto} holds, the model is both \emph{confident} ($\kappa_n$
low) and \emph{certain the fault is not internal} ($f_{int}$ low), so the
conventional trip is suppressed.
The $f_{thresh}$ is function-specific: $0.60$ for 87T and 87B; $0.50$
for 87L (tighter because channel-induced differential is weaker than
CT saturation).

\subsection{Layer 5 — Fallback and Coordination}
\label{sec:arch:l5}

Each function provides a fail-safe fallback when $\gamma < \gamma_{\min}$:
OC uses the greedy cascade of fixed commissioning settings~\cite{EluvathingalPaperB};
87T retains the conventional inrush/overexcitation blocking;
87L activates single-end OC backup in Loss mode;
87B retains the high-set instantaneous element.
All fallbacks are independent of the estimation outcome.

% =============================================================================
\section{SyncOC Function (M1--M4)}
\label{sec:syncoc}
% =============================================================================

The SAMBP-SyncOC framework is fully described in the companion
papers~\cite{EluvathingalPaperA,EluvathingalPaperB}; we summarise the
layer~2 and key results here for cross-function comparison.

\subsection{Zone Model}
\label{sec:syncoc:model}

The six-parameter post-fault current model
\begin{equation}
  i_A = I_{ss} + I_{sub}e^{-(t-t_0)/\tau_{ac}}\cos(\omega t+\phi_a)
        - I_{sub}\sin(\phi_a)\,e^{-(t-t_0)/\tau_{dc}}
  \label{eq:syncoc_model}
\end{equation}
uses the physics constraint $I_{dc} = -I_{sub}\sin\phi_a$, eliminating
the DC offset amplitude as a free parameter.
Without this constraint, a 7-parameter model yields $\kappa_n \approx
2\times10^8$; with it, $\kappa_n < 21$ for windows $\geq 67$\,ms.
The Hilbert--Fortescue extension (M3)~\cite{EluvathingalPaperB} handles
asymmetric faults via positive-sequence decomposition.

\subsection{Key Results}

Confidence scores $\gamma \in [0.89, 0.93]$ for 3PH faults;
$\gamma = 0.914$ for SLG (pickup reduced $-8.3\%$);
$\gamma = 0.587$ for LL (gate rejects, fixed settings retained).
The N-generator cascade (M4)~\cite{EluvathingalPaperB} enforces
selectivity for $N \in \{2,3,4\}$ in a single $\mathcal{O}(N)$ pass.

% =============================================================================
\section{Transformer Differential Function (87T)}
\label{sec:t87}
% =============================================================================

\subsection{Zone Model}
\label{sec:t87:model}

The 87T zone model uses five parameters
$\theta_T = [I_{diff,fund},\,\phi,\,k_2,\,k_5,\,\varepsilon_{CT}]$,
where $I_{diff,fund}$ is the fundamental differential current,
$k_2$, $k_5$ are the second- and fifth-harmonic ratios (inrush and
overexcitation indicators), and $\varepsilon_{CT}$ captures CT error.
The internal-fault probability $f_{int}$ is \emph{derived} as
$f_{int} = 1 - \max(k_2/k_{2,thresh},\,k_5/k_{5,thresh})$,
clipped to $[0,1]$.
This derivation is essential: at $\phi \approx 0$ (near-unity power factor
during transformer energisation), $f_{int}$ becomes collinear with
$I_{diff,fund}$, making it unidentifiable as a fitted parameter.

\subsection{Event Classification}
\label{sec:t87:results}

Fig.~\ref{fig:87T_char} shows the dual-slope differential characteristic
with event classification overlaid.
Table~\ref{tab:87T_events} summarises the 7 test events.
The critical result is the external fault + CT saturation event:
$\kappa_n = 4.95$, $f_{int} = 0.454 < 0.60$; the model veto
correctly suppresses the conventional trip, avoiding a false operation
that harmonic blocking alone would produce.
All 7/7 events are correctly classified.

\begin{figure}[!t]
  \centering
  \includegraphics[width=\columnwidth]{figures/87T_characteristic.png}
  \caption{87T differential characteristic with 7 test events overlaid.
           Green: no-trip region (normal load, blocked inrush/overexcitation).
           Blue: model-veto region (external fault + CT saturation).
           Red: trip region (three internal fault events).
           The dual-slope boundary ($I_{rst} = 1.0$, slopes 0.3/0.6\,pu)
           is from IEC~60255-151~\cite{IEC60255-151}.}
  \label{fig:87T_char}
\end{figure}

\begin{table}[!t]
  \renewcommand{\arraystretch}{1.15}
  \caption{87T event classification results (7/7 correct)}
  \label{tab:87T_events}
  \centering
  \small
  \begin{tabular}{lccccl}
    \toprule
    Event & $\kappa_n$ & $f_{int}$ & $\gamma$ & Result & Correct?\\
    \midrule
    Normal load      & 4.97 & 0.917 & 0.919 & no\_trip   & \checkmark \\
    Inrush           & 5.16 & 0.239 & 0.024 & blocked    & \checkmark \\
    Overexcitation   & 6.74 & 0.000 & 0.603 & blocked    & \checkmark \\
    Ext.+CT sat.     & 4.95 & 0.454 & 0.170 & model\_veto& \checkmark \\
    Int. A-G         & 4.83 & 0.865 & 0.323 & conv.~trip & \checkmark \\
    Int. A-B         & 4.83 & 0.847 & 0.011 & conv.~trip & \checkmark \\
    Int. 3PH         & 4.83 & 0.860 & 0.008 & conv.~trip & \checkmark \\
    \bottomrule
  \end{tabular}
\end{table}

% =============================================================================
\section{Line Current Differential Function (87L)}
\label{sec:l87}
% =============================================================================

\subsection{Zone Model and Channel Model}
\label{sec:l87:model}

The 87L zone model uses four parameters
$\theta_L = [I_{fund},\,\phi_{fund},\,I_{DC},\,\tau_{DC}]$,
where $I_{fund}$ is the fundamental differential current amplitude and
$I_{DC}$, $\tau_{DC}$ describe the DC exponential offset.
$f_{int}$ is derived as $\max(I_{fund}/I_{thresh},\,I_{DC}/I_{thresh})$,
clipped to $[0,1]$.
The DC exponential component is spectrally orthogonal to the sinusoidal
fundamental — unlike adding a third sinusoid at the same frequency (which
would be linearly dependent), $I_{DC}e^{-t/\tau_{DC}}$ contributes a
distinct column to the Jacobian, making all four parameters identifiable.
This is why a 4-parameter model is used rather than a 6-parameter one.

The channel model introduces three degradation modes:
\textit{Healthy} ($\varepsilon_{sync} = 0$, full 87L);
\textit{Degraded} ($\varepsilon_{sync} > 0$, sync-error veto active);
\textit{Loss} (heartbeat timeout, single-end OC backup).
Channel alignment uses circular roll (not zero-padding) to avoid
spurious differential at the buffer edge.

\subsection{Results Across 12 Scenarios}
\label{sec:l87:results}

Fig.~\ref{fig:87L_fsm} shows the three-mode FSM and the 12-scenario
classification matrix.
All healthy- and degraded-channel scenarios (8/8) are correctly classified.
For loss-channel scenarios, the single-end OC backup correctly trips on
external and 3PH faults (2/2); for the loss-channel AG internal fault
the system correctly applies no-trip by design (channel loss removes the
differential information needed to distinguish internal from external) — this
is the 1 expected-none case. Total: 11/12, with 1 expected-none.

Key quantitative result: $\kappa_n = 2.45$ for all scenarios — the
lowest in the SAMBP suite — reflecting the strong orthogonality between
the DC and sinusoidal components in the 4-parameter model.

The $\pi$-section charging compensation study (TR-13) reduces the mean
charging residual from 1.43\,pu to $0.94 \times 10^{-3}$\,pu, enabling
the minimum detectable fault level $\alpha_{50}$ to be reduced from 0.07
to 0.05 per unit for AG faults.
Voltage-based correction (TR-15) further reduces $\alpha_{50}$ to 0.03
(AG) and 0.02 (3PH).

\begin{figure}[!t]
  \centering
  \includegraphics[width=\columnwidth]{figures/fig_87L_fsm.pdf}
  \caption{87L three-mode FSM (top) and 12-scenario classification
           result matrix (bottom).
           H = Healthy channel; D = Degraded; L = Loss.
           Green check marks: correct classification;
           grey: expected-none (no-trip by architecture design).}
  \label{fig:87L_fsm}
\end{figure}

% =============================================================================
\section{Busbar Differential Function (87B)}
\label{sec:b87}
% =============================================================================

\subsection{Zone Model}
\label{sec:b87:model}

The 87B zone model uses only three parameters
$\theta_B = [I_{diff},\,\phi,\,\varepsilon_{CT}]$.
Unlike the transformer, a busbar has no magnetising branch, so inrush
and overexcitation harmonics are absent; $k_2$ and $k_5$ are excluded
from $\theta_B$ by physics.
$f_{int} = \mathrm{clip}(1 - \varepsilon_{CT}/0.10,\;0,\;1)$: when
CT error is large ($\varepsilon_{CT} \approx 0.10$), the differential
is externally driven and $f_{int} \to 0$; when CT error is negligible,
$f_{int} \to 1$.

\subsection{Event Classification}
\label{sec:b87:results}

Fig.~\ref{fig:87B_char} shows the 87B characteristic with six test events.

\begin{figure}[!t]
  \centering
  \includegraphics[width=\columnwidth]{figures/87B_characteristic.png}
  \caption{87B differential characteristic with 6 test events overlaid.
           Normal load and external fault lie in the no-trip region
           (conv\_trip = 0). External fault + CT saturation is intercepted
           by the model veto. Internal AG and 3PH faults trip
           conventionally. CT open-circuit is an architectural boundary
           case (see §\ref{sec:b87:limit}).}
  \label{fig:87B_char}
\end{figure}

Five of six events are correctly classified (Table~\ref{tab:87B_events}).
The CT saturation veto fires correctly: $\kappa_n = 2.73 < 30$,
$f_{int} = 0.00 < 0.60$.

\begin{table}[!t]
  \renewcommand{\arraystretch}{1.15}
  \caption{87B event classification results (5/6 correct)}
  \label{tab:87B_events}
  \centering
  \small
  \begin{tabular}{lccccl}
    \toprule
    Event & $\kappa_n$ & $f_{int}$ & $\gamma$ & Result & Correct?\\
    \midrule
    Normal load      & 4.36 & 1.000 & 0.000 & no\_trip   & \checkmark \\
    External fault   & 4.27 & 1.000 & 0.000 & no\_trip   & \checkmark \\
    Ext.+CT sat.     & 2.73 & 0.000 & 0.024 & model\_veto& \checkmark \\
    Internal A-G     & 4.27 & 1.000 & 0.137 & conv.~trip & \checkmark \\
    Internal 3PH     & 4.27 & 1.000 & 0.137 & conv.~trip & \checkmark \\
    CT open-circuit  & 4.36 & 1.000 & 0.845 & model~trip & $\times$\dag \\
    \bottomrule
    \multicolumn{6}{l}{\dag Known limitation: see §\ref{sec:b87:limit}.}
  \end{tabular}
\end{table}

\subsection{CT Open-Circuit Limitation}
\label{sec:b87:limit}

When a CT open-circuits, the primary current $I_{CT} \to 0$ and
$\varepsilon_{CT} \approx 0$.
The model assigns $f_{int} = 1$, the veto does not fire, and the model
issues a trip.
This is a known architectural boundary: the three-parameter zone model
cannot distinguish a CT open-circuit ($I_{diff}$ below threshold but
$\varepsilon_{CT} = 0$) from a genuine high-impedance internal fault.
A supplementary high-impedance (hi-Z) detection scheme is required for
full immunity, consistent with practice~\cite{Schweitzer2009}.

% =============================================================================
\section{Cross-Function Comparison and Design Principles}
\label{sec:comparison}
% =============================================================================

\subsection{Zone Model Comparison}
\label{sec:comparison:models}

Table~\ref{tab:model_compare} compares the four zone models.
Fig.~\ref{fig:kappa} shows the $\kappa_n$ distribution across all 28
simulation scenarios grouped by function.

\begin{table*}[!t]
  \renewcommand{\arraystretch}{1.25}
  \caption{SAMBP zone model comparison across four protection functions}
  \label{tab:model_compare}
  \centering
  \begin{tabular}{lccccl}
    \toprule
    Function & Parameters & Derived parameter & $\kappa_n$ range
             & $\kappa_n$ (7/6/5-param unconstrained) & Rationale \\
    \midrule
    SyncOC (OC)   & 6 & $I_{dc} = -I_{sub}\sin\phi_a$ & $4.5$--$21$
                  & $\approx 2\times10^8$ & DC offset from AC phase at fault instant \\
    Transformer (87T) & 5 & $f_{int} = 1 - \max(k_2,k_5)$ & $4.8$--$6.7$
                  & $\approx10^4$--$10^5$ & $f_{int}$ collinear with $I_{diff}$ at $\phi\approx0$ \\
    Line diff. (87L)  & 4 & $f_{int} = \max(I_{fund},I_{DC})/I_{thresh}$ & $2.45$
                  & $\approx10^3$--$10^4$ & 3 sinusoids at same frequency are rank-2 \\
    Busbar (87B)      & 3 & $f_{int} = 1 - \varepsilon_{CT}/0.10$ & $2.7$--$4.4$
                  & $\approx10^3$ & No transformer harmonics ($k_2,k_5$ absent) \\
    \bottomrule
  \end{tabular}
\end{table*}

\begin{figure}[!t]
  \centering
  \includegraphics[width=\columnwidth]{figures/fig_kappa_comparison.pdf}
  \caption{Column-normalised Jacobian condition number $\kappa_n$
           for all 28 simulation scenarios grouped by protection function.
           Red dashed line: $\kappa_{thresh} = 30$.
           All constrained models remain well below threshold.
           The unconstrained SyncOC 7-parameter model ($\kappa_n \approx
           2\times10^8$, shown as annotation) demonstrates the necessity
           of the physics-reduction step.}
  \label{fig:kappa}
\end{figure}

\subsection{Physics-Reduction Proposition}
\label{sec:comparison:prop}

\begin{proposition}[Physics-parameter reduction]
\label{prop:reduction}
Let $\theta \in \mathbb{R}^p$ be the parameter vector of a zone model
$\hat{i}(\theta,t)$, and let $\theta_k = g(\theta_{-k})$ be an exact
relation derivable from network physics (where $\theta_{-k}$ denotes all
parameters except $\theta_k$).
Then the column-normalised Jacobian of the constrained model
$\hat{i}(\theta_{-k},t)$ with $\theta_k$ substituted satisfies
$\kappa_n(\mathbf{J}_{p-1}) \ll \kappa_n(\mathbf{J}_p)$.
Empirically, the reduction exceeds one order of magnitude in all four
SAMBP functions (Table~\ref{tab:model_compare}).
\end{proposition}

\begin{proof}[Proof sketch]
The unconstrained Jacobian $\mathbf{J}_p$ has columns
$\partial\hat{i}/\partial\theta_j$.
When $\theta_k = g(\theta_{-k})$, the column
$\partial\hat{i}/\partial\theta_k$ becomes a linear combination of the
remaining columns (by the chain rule applied to $g$), making
$\mathbf{J}_p$ rank-deficient.
The condition number of a rank-deficient matrix is infinite; numerical
regularisation in LM produces $\kappa_n \gg 1$.
Substituting the physics constraint removes this column, restoring
the matrix to full column rank and eliminating the degenerate singular
value.
The magnitude of the reduction depends on the degree of collinearity;
for near-collinear cases (DC offset and AC at $t \approx t_0$), the
removal reduces $\kappa_n$ by six to eight orders of magnitude.
$\square$
\end{proof}

\begin{corollary}
For any protection zone model, an engineer should first check whether any
parameter can be expressed as a function of the others via network physics
(KVL/KCL, transformer turn-ratio, Fortescue identity).
If so, that parameter must be derived, not fitted, to achieve a
well-conditioned LM estimation.
\end{corollary}

\subsection{Stage-2 Model Veto — Cross-Function Characterisation}
\label{sec:comparison:veto}

Table~\ref{tab:veto_summary} summarises veto activation across all
four functions.
The veto fires in three of the 28 scenarios — one per 87T, 87L, 87B —
and in each case the physical cause is either CT saturation or channel
degradation.
In no case does the veto suppress a genuine internal fault.

\begin{table}[!t]
  \renewcommand{\arraystretch}{1.15}
  \caption{Model veto activations across all four functions}
  \label{tab:veto_summary}
  \centering
  \small
  \begin{tabular}{lcccc}
    \toprule
    Function & Veto fires? & $\kappa_n$ & $f_{int}$ & Physical cause \\
    \midrule
    SyncOC (OC) & No  & ---  & ---  & (no veto in OC path) \\
    87T  & Yes & 4.95 & 0.45 & External fault + CT sat. \\
    87L  & Yes & 2.45 & 0.00 & Degraded channel (healthy) \\
    87B  & Yes & 2.73 & 0.00 & External fault + CT sat. \\
    \bottomrule
  \end{tabular}
\end{table}

\subsection{Computational Cost}
\label{sec:comparison:compute}

Table~\ref{tab:compute} confirms that all four functions satisfy the
sub-cycle compute constraint.
The 87B function has the lowest cost (3 parameters, 80 NFev)
and the SyncOC function the highest (6 parameters, 200 NFev with
multi-generator cascade), but all remain below 20\,ms.

\begin{table}[!t]
  \renewcommand{\arraystretch}{1.15}
  \caption{Computational cost across all four functions}
  \label{tab:compute}
  \centering
  \begin{tabular}{lccc}
    \toprule
    Function & Parameters & NFev (typical) & Wall time (ms) \\
    \midrule
    SyncOC (M1--M4) & 6 & $\leq 200$ & $\leq 18$ \\
    87T             & 5 & $\leq 150$ & $\leq 12$ \\
    87L             & 4 & $\leq 100$ & $\leq 8$  \\
    87B             & 3 & $\leq 80$  & $\leq 6$  \\
    \bottomrule
  \end{tabular}
\end{table}

% =============================================================================
\section{Discussion}
\label{sec:disc}
% =============================================================================

\subsection{Known Limitations and Architectural Boundaries}
\label{sec:disc:limits}

Three failure modes are acknowledged as architectural boundaries:
\begin{itemize}
  \item \textit{SyncOC LL fault}: negative-sequence contamination inflates
        the LM residual, causing gate rejection~\cite{EluvathingalPaperB}.
        The framework falls back to fixed settings safely.
  \item \textit{87B CT open-circuit}: $\varepsilon_{CT} = 0$ gives
        $f_{int} = 1$; the veto cannot distinguish this from an internal
        fault. Supplementary hi-Z detection is required.
  \item \textit{87L channel-loss internal fault}: single-end OC cannot
        definitively identify internal vs.\ external faults without
        remote measurements. Expected no-trip is the safe failure mode.
\end{itemize}
None of these boundaries is a numerical deficiency; each has a physically
principled explanation and a defined safe fallback.

\subsection{Comparison with ML-Based Adaptive Protection}
\label{sec:disc:ml}

ML-based approaches~\cite{Sheikhzadeh2022,Gao2023} can achieve high
classification accuracy but require large labelled fault datasets,
offer limited physical interpretability, and have no established
path to certification under IEC~60255 or IEEE~C37.112.
The \textsc{sambp} framework is interpretable by design: every gate
decision traces to a physically meaningful parameter ($\kappa_n$,
$f_{int}$, $\gamma$), and the five-layer architecture preserves the
conventional relay as a safety net.
The trade-off is that \textsc{sambp} requires a physics-derivable
zone model for each new protection function — a constraint that is
straightforwardly satisfied for the four functions presented here.

\subsection{Hardware-in-the-Loop Validation Roadmap}
\label{sec:disc:hil}

Fig.~\ref{fig:hil} shows the planned M5 HIL validation timeline.
Software preparation (C port of the LM estimator,
IEC~61850~\cite{IEC61850} sampled-value interface) is ongoing.
RTDS hardware acquisition is planned for mid-2026~Q3;
HIL experiments across all four functions are targeted for 2026~Q4.
A revised Paper~C submission incorporating HIL results is planned
for early 2027.

\begin{figure}[!t]
  \centering
  \includegraphics[width=\columnwidth]{figures/fig_hil_roadmap.pdf}
  \caption{SAMBP M5 HIL validation roadmap.
           Green: complete. Blue: in progress.
           Orange: planned (software). Red: hardware-dependent.
           Today (April 2026) marker shown in red.}
  \label{fig:hil}
\end{figure}

% =============================================================================
\section{Conclusion}
\label{sec:conc}
% =============================================================================

\begin{enumerate}
  \item The \textsc{sambp} five-layer architecture unifies adaptive
        protection across OC, 87T, 87L, and 87B functions under a
        common design discipline, sharing the two-pass LM estimator
        and composite confidence gate.
  \item The Physics-Reduction Proposition (Proposition~\ref{prop:reduction})
        establishes that deriving one collinear parameter from network
        physics reduces $\kappa_n$ by at least one order of magnitude,
        enabling reliable LM convergence; this is demonstrated empirically
        across all four zone models ($\kappa_n < 30$ in all 28 scenarios).
  \item The Stage-2 model veto~\eqref{eq:veto} generalises CT saturation
        rejection (87T, 87B) and channel anomaly rejection (87L) as
        a single identifiability condition, correctly firing in 3/28
        scenarios without any false suppression of genuine internal faults.
  \item Combined classification accuracy is 26/28 correct;
        the 2 exceptions are a known architectural boundary (87B CT
        open-circuit) and a designed no-trip case (87L channel-loss).
  \item All four functions operate within 20\,ms on standard hardware.
        HIL validation on an RTDS platform is planned for 2026~Q4.
\end{enumerate}

% =============================================================================
\bibliography{references_c}
% =============================================================================

\end{document}
